\chapter{Lattice Impurity Model Classes}
\index{impurity model!lattice|textbf}

\label{ch:lattice-impurity-models}

\chapterstatus{proposed nested impurity-model hierarchy; no minimal model has
been selected.}

\section{Nested impurity model classes}
\index{model hierarchy!impurity}
\index{minimal model}


Stage~08 proposes a hierarchy that can distinguish scalar, orbital-dependent,
onsite, hopping, and range-restricted components. Let
$\mathfrak I_{0,d}\subseteq\mathfrak I_{1,d}\subseteq\cdots$ denote the frozen
nested family for dopant $d$. Each candidate is obtained by a declared
projection or fit of the accepted atomistic impurity operator.

A model level is evaluated with a vector of complementary errors,
\begin{equation}
  \boldsymbol\varepsilon_{m,d}
  =
  \left(
    \varepsilon_{H,m,d},
    \varepsilon_{\Pi,m,d},
    \Delta E_{b,m,d},
    1-F_{m,d}
  \right),
\end{equation}
where the components represent, respectively, an operator error, a retained
subspace or projector error, a bound-state energy error, and loss of a declared
fidelity measure. Each component still requires its exact definition, units,
normalization, and tolerance in the owning specification.

The selected impurity model is the least complex member satisfying every
prespecified tolerance. If no reduced member passes, the full atomistic
operator remains the accepted model. Failure to reduce is therefore a valid
scientific outcome rather than a software failure.
